\begin{abstract}
Reinforcement Learning from Human Feedback assumes that human annotators provide a stable preference signal --- yet the annotators themselves may adapt toward the AI style they repeatedly evaluate. We study this \emph{annotation drift} phenomenon in two real RLHF datasets (Anthropic HH-RLHF, 160,800 comparisons; OpenAI WebGPT, 19,578 comparisons) using the Alignment Asymmetry Index (AAI), a composite instrument that measures directional stylistic adaptation in preference labels. Our key finding is that verbosity preference reverses sign across annotation strata: early-round annotators penalize verbose responses while later-round annotators reward them ($\Delta\beta_L = +0.080$, non-overlapping 95\% bootstrap confidence intervals). A geometric projection of annotation preference gradients onto an AI-typicality embedding vector is significantly positive ($p = 2.05\times10^{-5}$) and discriminant-valid (placebo $p = 0.48$), confirming that the drift is specifically toward AI-typical text rather than random stylistic variation. Our null results --- untestable ambiguity moderation and subthreshold within-annotator effects --- precisely characterize the minimum data requirements for confirming individual annotator causal adaptation. These findings suggest that multi-round RLHF annotation datasets may encode temporal non-stationarity in the preference signal, motivating drift-aware quality auditing before reward model training.
\end{abstract}
